\chapter{Overview and package layout}
\label{ch:overview}

\section{Design principles}

Three commitments shape the package, and most of its structure follows
from them.

The first is that derivatives are never written by hand.  Every nuclear
gradient and second derivative is obtained by differentiating the
energy code path with JAX.  This is what makes the CASSCF, state-averaged
CASSCF, and NEVPT2 layers share one implementation of the physics: a
correction to the energy expression propagates to all of its derivatives
without a separate derivation.

The second is that the package is layered, and the layer boundaries are
tested.  The kernel cannot import the layers above it; the analysis
layer cannot import the kernel at all.  The contract is enforced in the
test suite (see Chapter~\ref{ch:architecture}), so a violation is a
failing test rather than a matter of review discipline.

The third is that convergence is reported, not assumed.  Optimizers
return the gradient norm at the point they stopped at, and the package
distinguishes ``the optimizer stopped improving'' from ``the point is
stationary.''  Callers are expected to carry that distinction into
whatever they claim about a result.

\section{Package layout}

\begin{lstlisting}[language=bash]
src/diffmrc/
  core/       numerical kernel: rotations, CASCI, CASSCF, NEVPT2,
              stationary points, the frozen source of the chain
  analysis/   post-processing that reads converged objects: entropies,
              mutual information, spin and state identification,
              orbital identity and frontier windows; NumPy only
  policy/     selection rules over analysis output: composition,
              entropy-based and ranked orbital selection, verification
  driver/     composition layer: assembles the kernel into the
              differentiable SCF -> CASSCF -> NEVPT2 pipelines,
              orbital-point convergence, frozen-geometry factories
  report/     result objects and rendering
  cli/        the command-line entry point
\end{lstlisting}

The dependency direction runs downward only: \texttt{core} at the
bottom, \texttt{cli} at the top.  A user script may import any layer;
the layers themselves may not import upward.  This is what allows the
analysis layer to be tested and used with nothing but NumPy installed.

\section{How the documents are organized}

This manual is the reference.  Part~I covers installation, the API
surface, and the protocol details that the code reports but does not
enforce.  Part~II carries the methods: the design of the differentiable
chain, the second-derivative construction, the verification records
against independent implementations, the benchmark results, and the
consolidated limitation list.  Part~III is for development: layer
contracts, the test architecture, and the reproduction steps for the
paper.

The companion \emph{Tutorial} is the on-ramp: installation, a first
CASSCF gradient, state-averaged workflows, active-space selection, and
a full benchmark walkthrough, each with runnable code.
